\section{Introduction}\label{sec:intro}

Cause--effect benchmarks drive method development in causal discovery. For two variables the reference is the T\"ubingen collection \citep{mooij2016distinguishing}, 108 pairs of real measurements of which 102 are univariate, and a method is summarised by its accuracy at orienting them; for graphs it is a small number of real networks, most often the protein-signalling data of \citet{sachs2005causal} with 17 consensus arcs. A new rule is accepted when its number is larger than the numbers of earlier rules in a table. This paper asks how much information such a table carries. It does so by running a panel of direction rules, under one protocol written before the runs, on seven real pair collections and on the graph benchmarks, and by asking two questions of the results: how large a difference the benchmarks can resolve, and why dependence-based rules, which use the structure that causal discovery is built on, so rarely beat much simpler rules on real data.

We use the word \emph{baseline} for rules that read the scale of regression errors or the marginal density of the variables: RECI \citep{blobaum2018cause}, IGCI \citep{daniusis2010inferring} and, since 2026, a parameter-free compression score \citep{stienstra2026oneruler}. The \emph{dependence-based} rules are those that test a model of the conditional distribution: additive-noise models \citep{hoyer2009nonlinear,peters2014identifiability}, location-scale models \citep{immer2023loci}, pairwise likelihood-ratio rules for non-Gaussian linear models \citep{hyvarinen2013pairwise}, and amortised or learned models \citep{lorch2022amortized}. On the T\"ubingen pairs all of them lie between $55\%$ and $67\%$ in our hands, with intervals of $\pm 10$ points; on chamber recordings with a known graph the ranking is almost the reverse.

\paragraph{Contributions.} (1) A same-protocol comparison of twelve pair rules (thirteen with one variant) on seven real collections (T\"ubingen; four Causal Chambers collections including three new ones; the Sachs pairs), with the stored decision of every rule on every pair, and a reproduction audit against published numbers (Section~\ref{sec:audit}). Reproducing the published accuracy of the compression baseline exactly required one undocumented rounding convention, which by itself moves the weighted headline by six points. (2) A resolution analysis: with 102 pairs a paired test has 80\% power only for a difference of about 19 points, and of the 78 pairwise comparisons among 13 rules two survive a Holm correction (Section~\ref{sec:noise}). (3) A multiverse of defensible reporting choices that changes ranks by up to ten places (Section~\ref{sec:multi}). (4) Evidence on mechanisms, each tested by prediction, by control or by manipulation (Section~\ref{sec:mech}): the real benchmarks are largely solvable from the marginal distributions alone, validated without leakage between training and test variables; additive-noise rules are near-perfect when the residual scale does not depend on the cause and reversed when it does, shown by a matched simulation with a marginal control and a factorial map; and an attempt to turn this account into a predictive model fails, which bounds the explanation (Section~\ref{sec:predict}). (5) An accounting of what does not work: every attempt to build a better rule during the study (Appendix~\ref{app:negative}) and the pre-registration log of all 42 experiments (Appendix~\ref{app:prereg}).

\paragraph{What this paper does not claim.} It does not claim that published methods misreport their results. We audited one published table in detail and reproduced it (Section~\ref{sec:audit}); for the others we do not know the protocol used and run the methods ourselves. It does not claim that dependence-based methods are worse in general: they are among the best rules in the regimes that fit their assumptions and the worst in others, and a benchmark that mixes regimes averages them. It does not claim that the mechanisms we identify are the only ones. The predictive test of Section~\ref{sec:predict} fails, which means that the account is incomplete.

\paragraph{Outline.} Section~\ref{sec:related} places the work. Section~\ref{sec:data} describes the collections and rules, Section~\ref{sec:protocol} the protocol and its provenance. Section~\ref{sec:coll} gives the cross-collection results. Sections~\ref{sec:noise}--\ref{sec:subsets} quantify resolution and reporting freedom. Section~\ref{sec:mech} tests mechanisms. Section~\ref{sec:graphs} reports the graph benchmarks. Section~\ref{sec:audit} is the reproduction audit. Section~\ref{sec:discussion} discusses and recommends; Section~\ref{sec:limits} lists the limits.
